% Generated draft: requires mathematical and source review.
\documentclass[11pt]{article}
\usepackage[margin=25mm]{geometry}
\usepackage{fontspec}
\setmainfont{Latin Modern Roman}
\usepackage{amsmath,amssymb,mathtools,unicode-math}
\setmathfont{Latin Modern Math}
\usepackage{enumitem,needspace,xurl,hyperref}
\setlength{\parindent}{0pt}
\setlength{\parskip}{5pt}
\setlength{\emergencystretch}{3em}
\newcommand{\N}{\mathbb{N}}
\newcommand{\Z}{\mathbb{Z}}
\newcommand{\Q}{\mathbb{Q}}
\newcommand{\R}{\mathbb{R}}
\newcommand{\C}{\mathbb{C}}
\newcommand{\F}{\mathbb{F}}
\newcommand{\A}{\mathbb{A}}
\newcommand{\PP}{\mathbb{P}}
\newcommand{\E}{\mathbb{E}}
\newcommand{\eps}{\varepsilon}
\newcommand{\dd}{\,\mathrm{d}}
\begin{document}
\section*{UnsolvedMath ID: 800462}
\textbf{Source problem code:} catalogue entry
\subsection{Definitions and mathematical statement}
In the usual Riemannian interpretation, let \((M^n,g,x)\) denote an \(n\)-dimensional Riemannian manifold with a distinguished point \(x\in M\). For \(q\in M\), let \(B_1(q)=\{y\in M:d_g(q,y)<1\}\), and let \(\mathrm{Vol}_g\) denote Riemannian volume. Interpret \(R\) as the scalar curvature of \(g\), and interpret the displayed Ricci bound as the tensor inequality
\[
\operatorname{Ric}_g(w,w)\geq-\lambda g(w,w)
\qquad\text{for every }q\in M\text{ and }w\in T_qM.
\]
These interpretations of \(R\) and the Ricci bound are conventional; the supplied statement does not explicitly define them. Whenever \(0<\mathrm{Vol}_g(B_1(x))<\infty\) and the integral is defined, write
\[
\fint_{B_1(x)}R\,\mathrm{d}\mathrm{Vol}_g
:=\frac{1}{\mathrm{Vol}_g(B_1(x))}
\int_{B_1(x)}R\,\mathrm{d}\mathrm{Vol}_g.
\]
The supplied assertion, attributed to Naber as Conjecture 6 in the context of noncollapsing Yau's conjecture, is that there is a positive constant \(C(n,\lambda,v)\), depending only on the indicated parameters, such that
\[
\operatorname{Ric}_g\geq-\lambda g,
\qquad \mathrm{Vol}_g(B_1(p))>v>0
\quad\Longrightarrow\quad
\fint_{B_1(x)}R\,\mathrm{d}\mathrm{Vol}_g<C(n,\lambda,v).
\]
The constant is intended to be uniform over the geometric objects satisfying the hypotheses. However, the supplied statement uses \(p\) in the volume hypothesis and \(x\) in the conclusion without specifying \(p\), its quantifier, or its relation to \(x\). In particular, the hypothesis has not been changed to \(\mathrm{Vol}_g(B_1(x))>v\). The admissible range of \(\lambda\), completeness, boundary conditions, and any restrictions on dimension are also unstated. These unresolved qualifications prevent a fully specified conjecture.
\subsection{Short English statement}
The supplied conjecture asks for a uniform upper bound on the average scalar curvature over a unit ball, depending only on dimension, a lower Ricci curvature bound, and a positive lower volume bound. The volume hypothesis concerns \(B_1(p)\), whereas the average concerns \(B_1(x)\); the relation between these centers and other geometric hypotheses requires clarification. The registry labels the problem open; this draft does not independently verify that status.
\subsection{Sources}
\begin{itemize}
\item Local UnsolvedMath registry, permanent ID \(800462\), source problem code \texttt{catalogue entry}, titled ``Spaces with Lower Ricci Curvature --- Conjecture 6 (noncollapsing Yau's conjecture).'' The supplied record attributes the problem to Naber, gives proposed year 2020, and labels its status as open.
\item Supplied source reference: Naber, \emph{Conjectures and Open Questions on Spaces with Lower Ricci Curvature} (2020), Conjecture 6, \url{https://arxiv.org/abs/2010.10031}. The record reports access on 2026-10-16 and extraction by \texttt{arxiv\_latex}. These are supplied provenance claims; the source has not been independently visited or verified for this draft.
\item Supplied attribution/context note: ``noncollapsing Yau's conjecture.'' The record supplies no separate original citation for this attribution.
\end{itemize}
\end{document}
